Um die Zylinderfläche als reelle Mannigfaltigkeit zu zeigen, betrachten wir die Konstruktion des Zylinders durch die Verklebung der Kanten eines Quadrats. Sei das Quadrat in der Ebene durch die Punkte \((0,0)\), \((1,0)\), \((1,1)\) und \((0,1)\) gegeben. Die Verklebung der vertikalen Kanten \((0,y)\) und \((1,y)\) für \(y \in [0,1]\) erfolgt durch die Identifikation \((0,y) \sim (1,y)\).

Ein Atlas für den Zylinder kann durch zwei Karten gegeben werden. Die erste Karte \((U_1, \varphi_1)\) sei definiert auf der offenen Menge \(U_1 = (0,1) \times (0,1)\) mit der Kartenabbildung \(\varphi_1: U_1 \to \mathbb{R}^2\), \(\varphi_1(x,y) = (x,y)\). Die zweite Karte \((U_2, \varphi_2)\) deckt die Umgebung der verklebten Kanten ab und sei definiert durch \(U_2 = (-\epsilon, 1+\epsilon) \times (0,1)\) für ein kleines \(\epsilon > 0\), mit der Kartenabbildung \(\varphi_2: U_2 \to \mathbb{R}^3\), \(\varphi_2(x,y) = (\cos(2\pi x), \sin(2\pi x), y)\).

Um die Glattheit der Übergangsfunktionen zu zeigen, betrachten wir den Schnitt \(U_1 \cap U_2 = (0,1) \times (0,1)\). Die Übergangsfunktion \(\varphi_2 \circ \varphi_1^{-1}\) ist gegeben durch:

\[
\varphi_2 \circ \varphi_1^{-1}(x,y) = \varphi_2(x,y) = (\cos(2\pi x), \sin(2\pi x), y)
\]

Da \(\cos(2\pi x)\) und \(\sin(2\pi x)\) glatte Funktionen sind, ist \(\varphi_2 \circ \varphi_1^{-1}\) glatt.

Um die Holomorphie der Übergangsfunktionen zu prüfen, beachten wir, dass die Übergangsfunktion \(\varphi_2 \circ \varphi_1^{-1}\) in der Form \((\cos(2\pi x), \sin(2\pi x), y)\) keine komplexen Komponenten enthält und somit trivialerweise holomorph ist, da sie konstant in der \(y\)-Komponente ist und die \(x\)-Abhängigkeit durch glatte trigonometrische Funktionen beschrieben wird.

Zusammenfassend ist die Zylinderfläche eine reelle Mannigfaltigkeit, und der angegebene Atlas definiert auch eine komplexe Struktur.

%CHECK: Adressiert die fehlende Spezifikation der Kartenabbildung \(\varphi_2\), die Glattheit der Übergangsfunktionen und die Holomorphie durch explizite Berechnungen.